\section{Background and related work}\label{sec:background}

\subsection{Shielded state transitions}

A shielded payment system commonly represents an unspent note by a commitment
inside an authenticated set and represents its consumption by a public,
unique nullifier.  A proof relates the nullifier and new commitment to a
private opening and membership path without revealing the spent note.
Zerocash gives an early complete treatment of this pattern in a decentralized
payment system \cite{BenSassonEtAl2014Zerocash}.  \Profile{} retains only the
cryptographic state-transition atom needed here: one input, one output, one
asset, a public fee, and a same-path replacement in a depth-20 private tree.
Wallet behaviour, note discovery, transaction-origin privacy, and multi-party
payment protocols are outside the relation.

It is useful to separate two statements that are sometimes conflated.  The
proof relation establishes that private note data are consistent with public
inputs.  The account transition establishes that the program consumed the
expected public pre-state and wrote the specified post-state.  The latter also
depends on the host runtime's account-ownership, rollback, and locking rules;
it is not a consequence of polynomial-commitment soundness.

\subsection{IOPs of proximity and transparent arguments}

An interactive oracle proof (IOP) lets a verifier query committed oracle
messages instead of reading them in full
\cite{BenSassonChiesaSpooner2016}.  The Fast Reed--Solomon IOP of Proximity
(FRI) tests
proximity to a Reed--Solomon code through recursive folding and a small number
of authenticated queries
\cite{BenSassonBentovHoreshRiabzev2018FRI}.  This approach underlies
transparent STARK systems
\cite{BenSassonBentovHoreshRiabzev2018STARK}.  DEEP-FRI strengthens the
analysis by sampling outside the evaluation domain
\cite{BenSassonGoldbergKoppartySaraf2020}.  WHIR combines multilinear-polynomial evaluation
with Reed--Solomon proximity testing and is designed for very fast
verification \cite{ArnonChiesaFenziYogev2025WHIR}.

The protocol in this paper belongs to that lineage but has its own transcript
and proof format.  It uses a circle-group evaluation domain, rate $1/512$, two
commitment layers, two out-of-domain/layer steps, four arity-four folds, and a
degree-four final polynomial.  Query positions address four-slot fibers.  The
transcript also includes a three-branch \(\mathsf{Good}_{\mathrm{spend}}\) selector and explicit
work nonces.  Consequently, WHIR's parameters, proof grammar, and theorem
statements are not imported by name; each lemma used here is specialized to
the custom schedule.

\subsection{Circle-domain Reed--Solomon codes}

The base field is $\Fbase$, often called M31.  Its multiplicative group does
not directly provide the large smooth domains used by conventional radix-two
fast Fourier transform (FFT) constructions.  Reed--Solomon codes over a
norm-one circle group provide
smooth evaluation domains in this setting
\cite{HaboeckLubarovNabaglo2023}.
Circle STARKs develop the corresponding arithmetization and proximity-testing
machinery for Mersenne-prime fields
\cite{HaboeckLevitPapini2024}.  The S-two
construction provides a closely related concrete circle-code transcript and
soundness development \cite{CarmonEtAl2026Stwo}.

\Profile{} uses the circle domain as a code domain while retaining a
multilinear claim at the protocol boundary.  The construction section defines
the resulting ordering, fiber map, extension-field arithmetic, and fold maps
explicitly.  This distinction matters because an isomorphism of domains alone
does not establish that a proximity theorem, batching argument, or query
bound transfers to a different transcript.

\subsection{Proximity gaps and correlated agreement}

Soundness for a folded proximity protocol requires more than the distance of
an individual Reed--Solomon word.  The analysis must control lists of nearby
codewords and show that correlated linear combinations do not allow
inconsistent candidates to survive successive folds.  Proximity-gap results
formalize the first issue
\cite{BenSassonCarmonIshaiKoppartySaraf2023}.  Mutual correlated
agreement (MCA) formalizes the second.  Recent results establish MCA
preservation for polynomial generators and give Reed--Solomon MCA theorems in
the Johnson regime
\cite{BordageChiesaGuanManzur2025,Haboeck2025MCA}.

The distinction between a proved Johnson-regime statement and a conjectural
capacity-regime statement is material.  Crites and Stewart give counterexamples
to Reed--Solomon proximity-gap conjectures used in earlier capacity-level
analyses \cite{CritesStewart2025}.  The \Profile{} parameter proof therefore does
not extrapolate a constant list bound to capacity.  It instantiates only the
proved proximity and MCA premises, then records every finite-field,
out-of-domain, folding, query, hash, and selector failure event in a single
ledger.  The exact specialization is given in
\Cref{sec:soundness,sec:hiding}.

\subsection{Fiat--Shamir compilation and hiding}

The implementation is non-interactive: verifier challenges and query indices
are derived from a domain-separated SHA-256 transcript using the
Fiat--Shamir transformation \cite{FiatShamir1987}.  Security of this
transformation is not captured by adding independent per-round errors.  The
BCS IOP framework gives a random-oracle compilation and treats private oracle
commitments \cite{BenSassonChiesaSpooner2016}; later work analyzes Fiat--Shamir security
for FRI and related SNARKs against bounded-query adversaries
\cite{BlockGarretaKatzThalerTiwariZajac2023}.  In \Profile{}, the interactive event union and
the random-oracle reduction are stated separately.  The latter counts 32
prover-message boundaries and normalizes false-acceptance probability by the
adversary's oracle-query work.

Hiding also depends on the complete transcript rather than only on the
arithmetic trace.  Full-field masks induce affine families of committed words,
private Merkle salts hide unopened leaf preimages in the declared oracle
model, and a simulator programs the transcript at fresh oracle inputs.  The
\(\mathsf{Good}_{\mathrm{spend}}\) condition is designed to make the public affine image and mask
fiber cardinality independent of the witness for every selectable schedule.
The security theorem includes the selector, variable proof length, work
nonces, serialization, program-visible logs, deterministic state mutation,
and the possibility that all bounded attempts abort.  This use of affine
masking and constrained codewords is related to zero-knowledge IOPPs for
constrained interleaved codes \cite{ChiesaFenziWeissenberg2026}; the concrete
\Profile{} simulator is specialized to its own fixed transcript and view.

\subsection{Hash functions and concrete arithmetic}

Arithmetic constraints and public note digests use Poseidon2 over M31;
Poseidon2 is a sponge-friendly permutation designed for efficient algebraic
implementations \cite{GrassiKhovratovichSchofnegger2023}.  Commitments, transcript
challenges, private-Merkle nodes, and work tests use domain-separated SHA-256.
These uses have different roles.  Poseidon2 collision resistance binds the
public arithmetic digests, while SHA-256 is modelled as the programmable random
oracle in the soundness and hiding games; SHA-256 itself is specified by
FIPS~180-4 \cite{NISTFIPS1804}.  The implementation pins the
concrete M31 Poseidon2 constants and arithmetic code to a specified Plonky3
revision \cite{Plonky3M31Poseidon2Version061}; the security section lists the corresponding
assumptions rather than treating a software dependency as a theorem.

\subsection{Solana execution and account state}

Solana transactions invoke programs over an explicit account list.  Accounts
carry an owner and arbitrary program data, and only the owning program may
modify that data under the runtime rules \cite{SolanaAccounts}.  Programs are
metered in compute units drawn from a transaction-wide budget
\cite{SolanaComputeBudget}.  A transaction either commits its account changes or
rolls them back on error \cite{SolanaTransactions,SolanaTransactionPipeline},
which makes verify-before-write control flow suitable for an atomic
proof-authorized transition.

Large proof transport is nevertheless distinct from verification.  \Profile{}
stores proof bytes in a program-owned account through a sequence of setup
transactions, then changes the account to an application-finalized state.  A
spend instruction refuses an unfinalized account and does not expose an upload
operation after finalization.  ``Finalized'' here describes this
application-level state machine; it is not a claim that account data are
intrinsically immutable at the protocol layer.  Likewise, a deployed SBF
program remains upgradeable while its loader-v3 upgrade authority is present
\cite{SolanaProgramDeployment}.  Deployment identity and authority state are
therefore recorded separately from the cryptographic proof identity.

Network observations use Solana's commitment levels.  In particular,
``finalized'' transaction evidence refers to a block recognized at the RPC
interface's finalized commitment level \cite{SolanaRPC}, which is
separate from the proof account's application-finalization bit.
